% Generated draft: requires mathematical and source review.
\documentclass[11pt]{article}
\usepackage[margin=25mm]{geometry}
\usepackage{fontspec}
\setmainfont{Latin Modern Roman}
\usepackage{amsmath,amssymb,mathtools,unicode-math}
\setmathfont{Latin Modern Math}
\usepackage{enumitem,needspace,xurl,hyperref}
\setlength{\parindent}{0pt}
\setlength{\parskip}{5pt}
\setlength{\emergencystretch}{3em}
\newcommand{\N}{\mathbb{N}}
\newcommand{\Z}{\mathbb{Z}}
\newcommand{\Q}{\mathbb{Q}}
\newcommand{\R}{\mathbb{R}}
\newcommand{\C}{\mathbb{C}}
\newcommand{\F}{\mathbb{F}}
\newcommand{\A}{\mathbb{A}}
\newcommand{\PP}{\mathbb{P}}
\newcommand{\E}{\mathbb{E}}
\newcommand{\eps}{\varepsilon}
\newcommand{\dd}{\,\mathrm{d}}
\begin{document}
\section*{UnsolvedMath ID: 30005876}
\textbf{Source problem code:} catalogue entry
\subsection{Definitions and mathematical statement}
For each place \(v\) of \(\Q\), let \(\Q_v\) denote the completion of \(\Q\); thus \(\Q_\infty=\R\), and the other completions are \(\Q_p\) for primes \(p\). For a variety \(X\) defined over \(\Q\), write \(X(K)\) for its set of \(K\)-rational points. Say that \(X\) satisfies the \emph{Hasse principle} (HP) if
\[
\left(X(\R)\ne\varnothing\ \text{and}\ X(\Q_p)\ne\varnothing\ \text{for every prime }p\right)
\Longrightarrow X(\Q)\ne\varnothing.
\]
There are two relevant meanings of unirationality. Unirationality over \(\Q\) means that, for some integer \(n\geq 0\), there is a dominant rational map
\[
\PP^n_{\Q}\dashrightarrow X
\]
defined over \(\Q\). Geometric unirationality means that, for some integer \(n\geq 0\), there is a dominant rational map
\[
\PP^n_{\overline{\Q}}\dashrightarrow X_{\overline{\Q}},
\qquad X_{\overline{\Q}}=X\times_{\Q}\overline{\Q}.
\]
\textbf{Source question, with unresolved scope.} The supplied record attributes to Colliot-Th\'el\`ene the conjecture that every unirational variety over \(\Q\) satisfies HP. Its quantified form is
\[
\forall X/\Q,\qquad
\left(\mathsf U(X)\ \land\ \forall v\;X(\Q_v)\ne\varnothing\right)
\Longrightarrow X(\Q)\ne\varnothing,
\]
where \(v\) ranges over all places of \(\Q\). Here \(\mathsf U\) records the source's unresolved meaning of unirationality: the excerpt does not distinguish unirationality over \(\Q\) from geometric unirationality. It also does not specify whether the varieties quantified over must be smooth, proper, projective, or geometrically integral. Consequently this is a faithful statement template requiring clarification, rather than an unqualified assertion of either interpretation.
\subsection{Short English statement}
The supplied question asks whether a unirational variety defined over \(\Q\) must have a rational point whenever it has a point over every completion of \(\Q\). The intended meaning of unirationality and the geometric hypotheses on the variety require clarification. The conjecture is attributed in the supplied excerpt to Colliot-Th\'el\`ene.
\subsection{Sources}
\begin{itemize}
\item Local registry: \texttt{data/unsolvedmath/problems.json}, permanent UnsolvedMath ID \(30005876\), source problem code \texttt{catalogue entry}, title \emph{The Hasse Principle for Unirational Varieties}. The supplied record labels the problem open and its statement \texttt{corrected\_verified}; these are supplied registry labels, not independent mathematical verification. Its verification note states that only formula delimiters, notation, grammar, and typography were repaired, with mathematical content and scope preserved.
\item Supplied report citation: \emph{Mini-Workshop: Growth and Expansion in Groups} (2024), Oberwolfach Reports, pp.\ 1005--1040, DOI \texttt{10.4171/owr/2024/17}. \url{https://doi.org/10.4171/owr/2024/17}. The excerpt attributes the conjecture to Colliot-Th\'el\`ene and includes the presentation title \emph{The group large sieve method and Hilbert's irreducibility theorem}, Daniele Garzoni, joint work with Lior Bary-Soroker. This presentation context does not identify those speakers as the conjecture's originators.
\item Supplied literature source, labelled ``current literature review'': \url{https://openscience.ub.uni-mainz.de/server/api/core/bitstreams/69271cb7-02cf-42fd-abc9-837e2bf69aa8/content}. No author, title, or pinpoint reference for this document is supplied.
\item Supplied literature assessment checked on \textbf{2026-08-22}: the record labels the general conjecture open and states that many classes of unirational varieties over \(\Q\) are known to satisfy the Hasse principle, while no theorem covers all such varieties. This is supplied dated curation derived from the OWR material, not an independently verified assessment. The supplied background explicitly requests independent verification of time-sensitive status and cited claims. Its listed sources are the literature-review link above and \url{https://doi.org/10.4171/owr/2024/17}.
\end{itemize}
\end{document}
